\chapter{最终结果和实验分析}

\section{端到端结果}
我们最终采用完整 BuildStorm 耗时作为主指标，因为短 pc-hot 只能告诉我们某个符号变热或变冷，不能证明比赛测例真的变快。已确认的空页池 A/B 结果如下：

\begin{table}[htbp]
  \centering
  \caption{预清零页面池的代表性 A/B 结果}
  \begin{tabular}{lrrr}
    \toprule
    配置 & 完整编译时间 & 加速比 & 耗时降低 \\
    \midrule
    优化前：前台同步清零 & 530\,s & 1.000$\times$ & -- \\
    优化后：idle 预清零页面池 & 502\,s & 1.056$\times$ & 5.28\% \\
    \bottomrule
  \end{tabular}
\end{table}

加速比按
\[
  S=\frac{T_{\mathrm{before}}}{T_{\mathrm{after}}}
\]
计算。这里的 530\,s 和 502\,s 是相同 workload 条件下的代表性完整轮结果；如果后续补充多轮交错 A/B，我们会以中位数替换单轮值，并保留原始日志路径。

\section{结果应该怎样解读}
这 28 秒不是由某一个短函数直接贡献的。页面清零被放到 idle CPU 后，前台 fault、页表建立和匿名映射路径少了一段同步工作；同时，批量从 allocator 取得 raw frame 也减少了全局分配器锁的获取次数。空页池优化因此和前面的 exec、mmap、page cache 改动存在路径上的叠加，但我们不把各项百分比简单相加，最终只报告完整组合和可复核的 A/B。

编译时间本身不是本轮所有优化的独立目标。我们主要优化内核运行时，因此构建时间数据只在同一宿主、同一工具链和同一配置下比较；不能把宿主编译缓存、QEMU 启动时间或不同镜像状态混入结论。功能回归、双架构 check 和 snapshot 运行是性能数据的前置条件，任何出现 panic、文件系统错误或测例缺项的候选都不进入最终组合。

\section{没有纳入最终组合的实验}
我们保留了失败实验的记录，但没有把它们写成成果。例如，process child index、u128 时间换算、路径规范化快速分支和 TLSF 懒统计都出现过局部指标改善；完整轮却没有稳定收益，甚至出现退化，所以最终回退。这个取舍很重要：我们愿意放弃一个局部漂亮的数字，换取完整 workload 上可重复的结果。
